\chapter{Synthetic Benchmark Generation and Evaluation Framework}
\label{ch:benchmarking}

\section{Motivation}
\label{sec:bench-motivation}

The fixed benchmark set supplied with the assignment --- fifteen valid
circuits and six intentionally invalid ones --- is well suited to
demonstrating correctness of the parser, the STA engine, and the
logical-effort analysis, and \cref{ch:results} reports results on it in
full. It is, however, a poor basis for comparing four different gate-sizing
heuristics against each other: fifteen circuits is too small a sample to
distinguish ``heuristic A is genuinely better than heuristic B'' from
``heuristic A happened to get a favorable draw of critical-path structure
on this particular handful of circuits.'' A statistically meaningful
comparison instead requires (i) a much larger population of test cases,
(ii) enough structural diversity across that population (circuit size,
depth, fan-out, reconvergence, and violation severity) to characterize
\emph{where} a heuristic wins or loses rather than only \emph{whether} it
does on average, and (iii) a known-good reference solution for each case
against which every heuristic's outcome can be measured, since without a
reference there is no way to distinguish ``heuristic A stopped early
because the problem was already solved'' from ``heuristic A stopped early
because it gave up.'' The benchmarking package addresses all three.

\section{Case Generation}
\label{sec:bench-generation}

\subsection{Topology Construction}

Two sources of circuit topology are supported: procedurally generated
levelized DAGs, constructed to a configured target size, depth, and
fan-out/reconvergence profile, and topologies imported unchanged from one of
the fixed valid seed circuits (\cref{ch:results}), allowing the same
generation-and-evaluation pipeline to be run against known, hand-inspectable
structures as a sanity check on the synthetic generator itself.

\subsection{Reference Solution via Multi-Start Beam Search}
\label{sec:bench-reference}

For a given topology, a \emph{best-known reference} gate-sizing assignment
is established independently of any of the four heuristics under test, by
a multi-start beam search that explores, from several different random
starting assignments, the space of legal per-gate size choices, retaining
at each search depth only the most promising partial assignments
(beam-limited to bound the search cost) and keeping the best complete
assignment found across all starts. This reference is explicitly certified
in the output only as \code{best\_known\_multi\_start\_beam} --- global
optimality is never claimed --- but it is a substantially stronger
reference than any single heuristic's own output, since it is not subject
to the same candidate-generation restrictions (top-5 pruning, single random
gate per iteration, and so on) that the heuristics themselves are.

\subsection{Planted Violations}
\label{sec:bench-planted}

Starting from the reference assignment, one or more gates are deliberately
\emph{downsized} until the resulting initial circuit exhibits a negative
WNS while remaining within its configured area and power budgets. This
produces a case with a known-repairable violation (the reference assignment
is a witness that a compliant, budget-respecting solution exists) and a
recorded ground truth for which specific gates were mutated away from the
reference (\code{planted\_mutations.csv}), which downstream evaluation
(\cref{sec:bench-evaluation}) can use as an additional diagnostic signal,
distinct from its primary correctness metric.

\subsection{Case Integrity and Reproducibility}
\label{sec:bench-integrity}

Both the initial (post-mutation) circuit and the reference assignment are
independently replayed through the production netlist parser, \code{Circuit}
constructor, and STA implementation --- the same code path used for
ordinary analysis --- before a case is considered valid and written to
disk, ensuring that no generated case can silently violate an invariant the
rest of the tool assumes. Every case directory is written atomically: if
generation is interrupted or exhausts its retry budget partway through, no
partial case directory is left behind, and suite-level failure diagnostics
(\code{generation\_failures.csv}) are preserved instead. All random seeds
used during generation are derived from SHA-256 hashes rather than from
Python's built-in, per-process-randomized \code{hash()} function, so that a
given configuration deterministically reproduces the same suite across
machines and Python invocations.

\subsection{Suite Layout}
\label{sec:bench-layout}

\begin{lstlisting}[language=bash, style=pythonstyle, numbers=none]
benchmarks/<suite>/
  suite_config.json          # the generating configuration
  cell_library.json          # copy of the library used for the suite
  suite_manifest.json        # index of every case and its parameters
  generation_cases.csv       # per-case generation record
  generation_failures.csv    # diagnostics for any discarded attempt
  cases/<case-id>/
    netlist.txt
    config.json
    reference_assignment.json
    benchmark_manifest.json
    planted_mutations.csv
\end{lstlisting}

\section{Evaluation Methodology}
\label{sec:bench-evaluation}

\subsection{Counterfactual Oracle}
\label{sec:bench-oracle}

The central evaluation instrument is an \emph{oracle} that is entirely
independent of any of the four heuristics being evaluated. For every gate
assignment (``state'') actually visited while replaying a heuristic's
recorded decisions on a case, the oracle evaluates \emph{every} legal
same-family cell replacement on \emph{every} gate in the circuit --- not
merely the gates a given heuristic happened to consider as candidates ---
and ranks them by the same two-phase rule used by the optimizer itself
(\cref{sec:opt-accept}): while a violation remains, rank by WNS
improvement, then TNS improvement on ties; once compliant, rank by cost
reduction subject to remaining compliant. This makes it possible to ask a
question no heuristic's own reported outcome can answer on its own: at each
individual decision point, was the choice the heuristic actually made the
\emph{best available} choice, a merely \emph{tied-best} choice, or a
strictly worse choice than something the oracle can show was available?
Exact recovery of the planted reference assignment
(\cref{sec:bench-planted}) is retained only as a secondary diagnostic,
since more than one gate-sizing assignment can be simultaneously compliant
and cost-optimal; the oracle-ranked decision quality at each step is the
primary correctness signal.

\subsection{Parallel, Deterministic Evaluation}
\label{sec:bench-parallel}

Cases are independent of one another and are evaluated in parallel (four
worker processes by default, configurable), with per-run timeouts enforced
so that a single pathological case cannot stall an entire suite; individual
case report fragments are merged back together in a fixed manifest order
so that the final evaluation report is reproducible regardless of which
worker happened to finish a given case first. Random Greedy, being
inherently stochastic, is additionally repeated with several deterministic
seeds per case so that its reported metrics reflect a distribution of
outcomes rather than a single draw.

\subsection{Reported Metrics}
\label{sec:bench-metrics}

\begin{longtable}{@{}p{0.35\textwidth}p{0.65\textwidth}@{}}
\toprule
\textbf{Artifact} & \textbf{Contents} \\
\midrule
\endhead
\code{case\_runs.csv} &
  Final metrics per case/heuristic: constraint-repair status, runtime,
  iteration and STA-call counts, and distance from the reference
  assignment. \\[4pt]
\code{optimizer\_iterations.csv} &
  A fully replayable, step-by-step history of every heuristic run. \\[4pt]
\code{oracle\_states.csv}, \code{oracle\_candidates.csv} &
  Every unique gate-assignment state visited, and every counterfactual
  candidate the oracle evaluated at that state, with feasibility,
  benefit, and tied-best status. \\[4pt]
\code{gate\_selection\_scores.csv} &
  Per gate/move: selection hit rate, WNS/TNS/cost regret relative to the
  oracle's best available choice, and missed off-path opportunities. \\[4pt]
\code{heuristic\_summary.csv} &
  \textbf{Wilson confidence intervals} on repair rate, failure and timeout
  rates, Random Greedy's runtime dispersion across its repeated seeds,
  selection accuracy against the oracle, and aggregate regret. \\[4pt]
\code{evaluation\_summary.json} &
  Hierarchical aggregates stratified by circuit source, size, depth,
  fan-out, reconvergence, violation severity, and slack headroom, plus a
  manifest of every artifact produced. \\
\bottomrule
\caption{Artifacts produced by \code{vlsi-sta benchmark evaluate}.}
\label{tab:bench-artifacts}
\end{longtable}

Using a \textbf{Wilson score interval} rather than a naive
$\hat p \pm 1.96\sqrt{\hat p(1-\hat p)/n}$ normal-approximation interval for
the repair-rate confidence bound is a deliberately more careful choice: the
Wilson interval remains well-behaved (bounded within $[0,1]$ and reasonably
calibrated) even when the observed repair rate is close to $0$ or $1$ or
the sample count for a given stratum is small, both of which are expected
to occur in some of the finer-grained strata of
\cref{sec:bench-metrics}'s hierarchical aggregation.

\section{Plotting}
\label{sec:bench-plotting}

\code{vlsi-sta benchmark plot} consumes an evaluation directory and
produces the following plots using the files listed on \Cref{tab:bench-artifacts}:

\begin{lstlisting}[numbers=none]
benchmarks/<suite>/evaluation/<run-id>/plots
├── case_reliability.png
├── heuristic_overview.png
├── repair_convergence.png
├── runtime_scalability.png
└── selection_quality.png
\end{lstlisting}

% \begin{figure}[H]
%     \centering
%     \includegraphics[width=0.7\textwidth]{benchmarking/heuristic_overview}
%     \includegraphics[width=0.7\textwidth]{benchmarking/runtime_scalability}
%     \includegraphics[width=0.7\textwidth]{benchmarking/repair_convergence}
%     \includegraphics[width=0.7\textwidth]{benchmarking/case_reliability}
%     \caption{Example output of the benchmark plotting command. Replace
%     with an actual generated plot from a completed evaluation run.}
%     \label{fig:bench-plot-placeholder}
% \end{figure}
